% Part: first-order-logic
% Chapter: proof-systems
% Section: natural-deduction

\documentclass[../../../include/open-logic-section]{subfiles}

\begin{document}

\iftag{FOL}
      {\olfileid{fol}{prf}{ntd}}
      {\olfileid{pl}{prf}{ntd}}

\olsection{సహజ నిగమనం}

సహజ నిగమనం, వాస్తవ తార్కిక ఆలోచనను (ముఖ్యంగా గణితశాస్త్రజ్ఞులు
ఉపయోగించే క్రమబద్ధమైన తార్కిక ఆలోచనను) ప్రతిబింబించడానికి ఉద్దేశించిన
ఒక !!a{derivation} వ్యవస్థ. వాస్తవ తార్కిక ఆలోచన అనేక “సహజ”
నమూనాల్లో సాగుతుంది. ఉదాహరణకు, సందర్భాల ద్వారా నిరూపణలో ఒక
వికల్ప పూర్వపక్షం ఆధారంగా నిష్కర్షను స్థాపించడానికి, ఆ వికల్పంలోని
రెండు వికల్పాంగాల్లో దేని నుంచైనా నిష్కర్ష వస్తుందని స్థాపిస్తాం.
పరోక్ష నిరూపణలో ఒక నిష్కర్ష నిషేధం వైరుధ్యానికి దారితీస్తుందని
చూపి ఆ నిష్కర్షను స్థాపిస్తాం. షరతీయ నిరూపణలో పూర్వాంగం నుంచి
ఉత్తరాంగం వస్తుందని చూపి “ఒకవేళ \dots అయితే \dots” అనే షరతీయ
ప్రకటనను స్థాపిస్తాం. ఈ సహజ నిగమనాల్లో కొన్నింటి రూపీకరణే సహజ
నిగమనం. ప్రతి తార్కిక సంయోజకానికీ, పరిమాణీకరణికకూ రెండు నియమాలు,
ఒక ప్రవేశ నియమం, ఒక తొలగింపు నియమం ఉంటాయి; వాటిలో ప్రతి ఒక్కటీ
అలాంటి ఒక సహజ నిగమన నమూనాకు అనుగుణంగా ఉంటుంది. ఉదాహరణకు,
$\Intro{\lif}$ షరతీయ నిరూపణకు, $\Elim{\lor}$ సందర్భాల ద్వారా
నిరూపణకు అనుగుణంగా ఉంటాయి. $!A \land !B$ నుంచి $!A$ను (లేదా
$!B$ను) నిగమించనిచ్చే $\Elim{\land}$ ఎంతో సరళమైన నియమం.

ఇతర !!{derivation} వ్యవస్థల నుంచి సహజ నిగమనాన్ని వేరు చేసే ఒక
లక్షణం పరికల్పనలను ఉపయోగించడం. సహజ నిగమనంలో ఒక !!^a{derivation}
అంటే !!{formula}s వృక్షం. !!{formula}s వృక్షపు మూలం వద్ద ఒకే
!!{formula} ఉంటుంది; వృక్షపు “ఆకులు” నిష్కర్షను నిగమించే
!!{formula}s. సహజ నిగమనంలో కొన్ని ఆకు !!{formula}s
!!{derivation} లోపల ఒక పాత్ర పోషించినప్పటికీ, !!{derivation}
నిష్కర్షకు చేరేసరికి “వాడుకొని ముగిసిపోతాయి.” ఇది వాస్తవ తార్కిక
ఆలోచనలో కొంతసేపు మాత్రమే అమల్లో ఉండే పరికల్పనలను ప్రవేశపెట్టే
ఆచారానికి అనుగుణం. ఉదాహరణకు, సందర్భాల ద్వారా నిరూపణలో ప్రతి
వికల్పాంగం సత్యమని పరికల్పిస్తాం; షరతీయ నిరూపణలో పూర్వాంగం
సత్యమని పరికల్పిస్తాం; పరోక్ష నిరూపణలో నిష్కర్ష నిషేధం సత్యమని
పరికల్పిస్తాం. ఊహాత్మక పరికల్పనలను ఈ విధంగా ప్రవేశపెట్టి, ఒక
మధ్యంతర మెట్టును స్థాపించే క్రమంలో తరువాత వాటిని తొలగించడం సహజ
నిగమనపు ప్రత్యేక లక్షణం. సహజ నిగమన !!{derivation}లోని ఆకుల వద్ద
ఉన్న సూత్రాలను పరికల్పనలు అంటారు; కొన్ని నిగమన నియమాలు వాటిని
“!!{discharge}” చేయవచ్చు. ఉదాహరణకు, $!A$ను కలిగి ఉన్న కొన్ని
పరికల్పనల నుంచి $!B$కు ఒక !!a{derivation} ఉంటే, $!A \lif !B$ను
నిగమించి $!A$ రూపంలోని ఏ పరికల్పననైనా ఉపసంహరించడానికి
$\Intro{\lif}$ నియమం అనుమతిస్తుంది. (ఏ నిగమనాల వద్ద ఏ పరికల్పనలు
ఉపసంహరించబడుతున్నాయో గుర్తుంచుకోవడానికి, ఆ నిగమనానికీ అది
ఉపసంహరించే పరికల్పనలకూ ఒక సంఖ్యతో గుర్తు పెడతాం.)
!!{derivation} చివరికి !!{undischarged}గా మిగిలే పరికల్పనలు కలిసి
నిష్కర్ష సత్యానికి సరిపోతాయి; కాబట్టి ఒక !!a{derivation}, తన
!!{undischarged} పరికల్పనలు తన నిష్కర్షను అనుగమిస్తాయని స్థాపిస్తుంది.

$!A$ వృక్షంలోని చివరి !!{sentence}గా ఉండి, !!{undischarged}గా ఉన్న
ప్రతి ఆకు $\Gamma$లో ఉండే ఒక !!a{derivation} ఉంటే, అప్పుడు మాత్రమే
సహజ నిగమనం ఆధారమైన $\Gamma \Proves !A$ సంబంధం ఉంటుంది. $!A$ చివరి
!!{sentence}గా ఉండి, అన్ని పరికల్పనలూ !!{discharged} అయిన ఒక
!!a{derivation} ఉంటే, అప్పుడు మాత్రమే $!A$ సహజ నిగమనంలో ఒక
సిద్ధాంతం. ఉదాహరణకు, $\Proves (!A \land !B) \lif !A$ అని చూపే ఒక
!!a{derivation} ఇది:
\begin{prooftree}
  \AxiomC{$\Discharge{!A \land !B}{1}$}
  \RightLabel{\Elim{\land}}
  \UnaryInfC{$!A$}
  \DischargeRule{\Intro{\lif}}{1}
  \UnaryInfC{$(!A \land !B) \lif !A$}
\end{prooftree}
$!A \land !B$ అనే పరికల్పన \Intro{\lif} నిగమనం వద్ద
!!{discharged} అయిందని $1$ గుర్తు సూచిస్తుంది.

సహజ నిగమనంలో $\Gamma \Proves \lfalse$ అయితే, అప్పుడు మాత్రమే
సమితి~$\Gamma$ వైరుధ్యభరితమైనది. వైరుధ్యభరిత సమితి నుంచి ఏ
!!{sentence}నైనా !!{derive} చేయగలిగేలా \FalseInt{} నియమం చేస్తుంది.

సహజ నిగమన వ్యవస్థలను 1930లలో గెర్హార్డ్ గెంట్సెన్, స్టానిస్వావ్
యాష్కోవ్స్కీ అభివృద్ధి చేశారు; తరువాత డాగ్ ప్రావిట్జ్, ఫ్రెడరిక్
ఫిచ్ మరింత అభివృద్ధి చేశారు. దాని నిగమనాలు సహజ నిరూపణ పద్ధతులను
ప్రతిబింబిస్తాయి కనుక తత్వవేత్తలు దానిని ఇష్టపడతారు. ఫిచ్
అభివృద్ధి చేసిన రూపాలు పరిచయాత్మక తర్కశాస్త్ర పాఠ్యపుస్తకాల్లో
తరచుగా ఉపయోగిస్తారు. తర్క తత్వశాస్త్రంలో, సహజ నిగమన నియమాలే
తార్కిక సంయోజకాల అర్థాలను ఇస్తాయని కొన్నిసార్లు భావించారు
(“నిరూపణ-సిద్ధాంత అర్థవిచారం”).

\end{document}
